\chapter{Propozycje scenariuszy ataków socjotechnicznych}\label{ch:socjotechnika}

\section{Założenia ogólne i ramy bezpieczeństwa projektowanych ataków}\label{sec:zalozenia-atakow}

\section{Scenariusz ataku nr 1}\label{sec:scenariusz-1}
\subsection{Cel, grupa docelowa i uzasadnienie na podstawie OSINT}\label{subsec:scen1-cel}
\subsection{Szczegółowy opis techniczny (narzędzia, payload, infrastruktura)}\label{subsec:scen1-techniczny}
\subsection{Mechanizmy ochronne i zabezpieczenia przed dostępem osób postronnych}\label{subsec:scen1-zabezpieczenia}
\subsection{Analiza ryzyka i procedury bezpieczeństwa gwarantujące brak szkód}\label{subsec:scen1-ryzyko}

\section{Scenariusz ataku nr 2}\label{sec:scenariusz-2}
\subsection{Cel, grupa docelowa i uzasadnienie na podstawie OSINT}\label{subsec:scen2-cel}
\subsection{Szczegółowy opis techniczny (narzędzia, payload, infrastruktura)}\label{subsec:scen2-techniczny}
\subsection{Mechanizmy ochronne i zabezpieczenia przed dostępem osób postronnych}\label{subsec:scen2-zabezpieczenia}
\subsection{Analiza ryzyka i procedury bezpieczeństwa gwarantujące brak szkód}\label{subsec:scen2-ryzyko}

\section{Scenariusz ataku nr 3}\label{sec:scenariusz-3}
\subsection{Cel, grupa docelowa i uzasadnienie na podstawie OSINT}\label{subsec:scen3-cel}
\subsection{Szczegółowy opis techniczny (narzędzia, payload, infrastruktura)}\label{subsec:scen3-techniczny}
\subsection{Mechanizmy ochronne i zabezpieczenia przed dostępem osób postronnych}\label{subsec:scen3-zabezpieczenia}
\subsection{Analiza ryzyka i procedury bezpieczeństwa gwarantujące brak szkód}\label{subsec:scen3-ryzyko}

